\section{Dataset-construction prompts}
The following prompts retain the complete system and user messages from the source implementations. The screening threshold is applied after scoring: all applicable consistency scores must be at least 7. Score intervals within a prompt are scoring rubrics, not additional retention thresholds.

\paragraph{D1: Viewpoint screening and change classification (Gemini-3-Flash-Preview).}
\begingroup\fontsize{8}{10}\selectfont
\begin{promptbox}
\begin{verbatim}
# Role You are a professional urban planner and computer vision expert. Your task
is to compare two panoramic street-view images (Image 1 and Image 2) of the same
location from different years and accurately identify real changes in the urban
physical environment.

# Task Execution Logic Silently perform the following two stages of analysis
internally. Do not output any analysis; output only one letter representing the
conclusion at the end.

## Stage 1: Strict viewpoint-alignment screening (mandatory prerequisite) Before
comparing changes, identify large static reference objects in Image 1 (such as
distant buildings, intersection vanishing points, and continuous building
outlines) and locate them in Image
2. Normal parallax (alignment successful): perspective deformation or slight
displacement caused by different capture lanes or a few meters of
forward/backward movement. Reference objects remain in approximately the same
regions. Continue to Stage 2. Major deviation (alignment failed): failed
north-heading correction produces a rotation of tens of degrees. For example, a
building at the exact center of Image 1 moves to the far left/right in Image 2,
or a view facing a building becomes a view facing a straight road. Stop
immediately and output E.

## Stage 2: Change detection (only after successful alignment) If alignment is
adequate, strictly distinguish changes using the following criteria.

Disturbances that MUST be ignored (if only these changes occur, classify as no
change):
1. Climate and lighting: seasonal tree appearance, standing water, cloud cover,
shadow position, and lighting angle.
2. Temporary/mobile objects: pedestrians, vehicles, construction barriers, and
temporary litter.
3. Small ancillary facilities: added or removed streetlights, traffic lights,
signs, fire hydrants, and trash bins.
4. Vegetation and barriers: added or removed trees and bushes; newly built or
removed walls and guardrails. Be extremely careful: never mistake a building
obscured by vegetation, a wall, or a guardrail for an actual change to the
building itself.

Core targets (count only clear changes): Buildings: complete demolition, new
construction, or major renovation of the main structure/exterior facade (such as
adding a second floor to a single-storey building). Ignore minor exterior-wall
color differences. Roads:
1. Major marking updates: large newly added crosswalks, directional arrows, or
road-surface text. Ignore additions, removals, or color changes of ordinary
dashed/solid road lines.
2. Clear material updates: dirt roads converted to asphalt, or conspicuous large
new pavement patches. Ignore minor crack changes, standing water, and reflections
caused by lighting.
3. Road-structure changes: new ramps or branches, or complete road removal.

# Output Format Strictly output exactly ONE uppercase English letter. Do not
include punctuation, spaces, line breaks, or explanation. Choose one of the
following five letters:
E: viewpoint alignment failed with a major deviation.
A: viewpoints aligned, with building changes only.
B: viewpoints aligned, with road changes only.
C: viewpoints aligned, with substantial changes to both buildings and roads.
F: viewpoints aligned, with no changes to the core targets.
\end{verbatim}
\end{promptbox}
\endgroup

The source labels E, A, B, C, and F map to viewpoint mismatch, building change, road change, building-and-road change, and no change, respectively.

\paragraph{D2: Cross-view consistency for changed samples (GPT-5.6-Luna).}
\noindent\textbf{System prompt.}\par
\begingroup\fontsize{8}{10}\selectfont
\begin{promptbox}
\begin{verbatim}
You are a geospatial consistency auditor evaluating the temporal and structural
alignment between top-down satellite orthophotos and 360-degree street-level
panoramas. Your baseline assumption is that discrepancies might exist due to time
mismatches; you must actively hunt for them.

CRITICAL RULE: The street-view images have been pre-rotated so the camera always
faces due North. Vertical reference lines and text labels are explicitly drawn on
the street-view images to guide you:
- Center = North (N): Marked by a thick solid white line and yellow 'N' text.
- Left-quarter = West (W): Marked by a dashed gray line and 'W' text.
- Right-quarter = East (E): Marked by a dashed gray line and 'E' text.
- Both left and right edges = South (S): Marked by dashed blue lines and 'S'
text.

SPATIAL DIRECTION MAPPING: To compare these two different perspectives, you must
understand how their directions align relative to the camera position (the RED
DOT in the satellite image):
- North (N): The center of the street-view corresponds to looking Up (Top) from
the red dot in the satellite image.
- East (E): The right-quarter of the street-view corresponds to looking Right
from the red dot in the satellite image.
- South (S): Both the far-left and far-right edges of the street-view wrap around
to correspond to looking Down (Bottom) from the red dot in the satellite image.
- West (W): The left-quarter of the street-view corresponds to looking Left from
the red dot in the satellite image.

STRICT AUDIT CRITERIA:
1. Blurry Image Rule (FATAL): If the street-view image contains any blurred areas
(whether small patches or large regions), you MUST immediately assign a
[Consistency Score] of 0. Blurriness makes the data unusable for high-quality
training.
2. Building Topology: Check for the presence, absence, relative height, and
footprint shape of permanent buildings in each direction.
3. Land Cover & Developmental State: Evaluate macro-level surface attributes
(e.g., dirt lot vs. paved road, active construction vs. finished building).
4. The "Temporary Object" Rule: Ignore individual mobile objects (e.g., vehicles,
pedestrians, shadows). However, you MUST evaluate permanent infrastructure.
5. Visibility Constraints: Ignore flat rooftop details in the satellite image
that cannot possibly be seen from the ground level.
\end{verbatim}
\end{promptbox}
\endgroup

\noindent\textbf{User prompt.}\par
\begingroup\fontsize{8}{10}\selectfont
\begin{promptbox}
\begin{verbatim}
Analyze the consistency between this street-view and satellite image pair based
on the spatial directional mapping and strict audit criteria. Consider all four
directions (North, East, South, West) in your analysis.

SCORING RUBRIC (0-10):
- 0: [FATAL] The street-view image is blurry (small or large areas), or the
locations are completely disjointed with massive contradictions.
- 1-3: Major mismatches in multiple directions.
- 4-6: Significant temporal changes in one or more directions (e.g., building
under construction vs. completed).
- 7-9: Minor structural differences but overall geometry and land use remain the
same.
- 10: Perfect match. All buildings, road layouts, and macro-level land uses match
perfectly across all directions.

You MUST output ONLY the score on the first line, in this exact format:

[Consistency Score]: <An integer from 0 to 10>
\end{verbatim}
\end{promptbox}
\endgroup

\paragraph{D3: Cross-View and Temporal Consistency Screening for No-Change Samples (GPT-5.6-Luna).}
The no-change branch contains two calls: a four-image audit of cross-view and satellite temporal consistency, and an earlier--recent street-view temporal audit. Their system/user prompt pairs are retained separately below; the four resulting scores must each be at least 7.
\subparagraph{D3(a): Cross-view and satellite temporal audit.}
\noindent\textbf{System prompt.}\par
\begingroup\fontsize{8}{10}\selectfont
\begin{promptbox}
\begin{verbatim}
You are an elite geospatial consistency auditor. You will be provided with ONE
combined image containing FOUR sub-images arranged in a 2x2 grid of the EXACT
SAME geographic location:

TOP-LEFT:    T1 Historical Street View (T1 Street View) — 2:1 panorama with
N/S/W/E direction markers. TOP-RIGHT:   T1 Historical Satellite (T1 Satellite) —
1:1 top-down orthophoto with red dot + yellow north arrow. BOTTOM-LEFT: T2
Current Street View (T2 Street View) — same format as top-left, current time.
BOTTOM-RIGHT: T2 Current Satellite (T2 Satellite) — same format as top-right,
current time.

*NOTE: Each quadrant has its title printed on the top-left corner in a black box.
Quadrants are separated by gray crosshair lines.*

================================================================== PART 1 --
STREET VIEW VISUAL GUIDES (top-left & bottom-left)
================================================================== The
street-view sub-images are 360-degree panoramas flattened into a 2D canvas,
pre-rotated so the camera always faces due North. Vertical reference lines and
text labels:
- Center (thick solid WHITE line + yellow 'N' text) = North
- Left-quarter (dashed GRAY line + 'W' text) = West
- Right-quarter (dashed GRAY line + 'E' text) = East
- Both left and right edges (dashed BLUE lines + 'S' text) = South

================================================================== PART 2 --
SATELLITE VISUAL GUIDES (top-right & bottom-right)
================================================================== Both satellite
sub-images have:
- A RED DOT in the center (exact same camera/focal location).
- A YELLOW ARROW pointing Up (due North).
- N, S, W, E text labels at the top, bottom, left, and right edges.

================================================================== PART 3 --
SPATIAL DIRECTION MAPPING
==================================================================
- North (N): Center of street-view = UP (TOP) from the red dot.
- East (E):  Right-quarter of street-view = RIGHT from the red dot.
- South (S): Far-left and far-right edges wrap = DOWN (BOTTOM) from the red dot.
- West (W):  Left-quarter of street-view = LEFT from the red dot.

================================================================== PART 4 -- YOUR
THREE TASKS ==================================================================

--- TASK A: T1 Cross-View Consistency (top-left vs top-right) --- Evaluate
whether T1 street-view (top-left) matches T1 satellite (top-right). Criteria:

A1. [FATAL BLUR RULE]: If the top-left street-view contains ANY blurred areas,
score MUST be 0. A2. Building Topology: Compare presence, absence, height,
footprint of permanent buildings in all four directions. A3. Land Cover &
Developmental State: dirt lot vs. paved road, construction vs. finished building,
etc. A4. "Temporary Object" Rule: IGNORE vehicles, pedestrians, shadows. DO
evaluate permanent infrastructure. A5. Visibility Constraints: IGNORE flat
rooftop details not visible from ground level.

--- TASK B: T2 Cross-View Consistency (bottom-left vs bottom-right) --- Evaluate
whether T2 street-view (bottom-left) matches T2 satellite (bottom-right). Apply
SAME criteria A1-A5.

--- TASK C: Satellite Temporal Consistency (top-right vs bottom-right) ---
Compare T1 satellite (top-right) vs T2 satellite (bottom-right). Verify
"NO-CHANGE" in physical urban structures.

C1. Focus on Static Structures (CRITICAL): Compare buildings, paved roads, tree
canopies in all four directions from red dot. C2. EXCLUSIONS (MUST IGNORE):
- Image resolution, blurriness, quality differences.
- Color, brightness, contrast, sensor changes.
- Seasonal changes (green vs. brown leaves, snow, shadows).
- Temporary/mobile objects (cars, tents, construction vehicles not altering
permanent structures). C3. "Perfect Match" Rule: If buildings and roads remain
identical, score is 10/10 regardless of color/season/clarity changes.
\end{verbatim}
\end{promptbox}
\endgroup

\noindent\textbf{User prompt.}\par
\begingroup\fontsize{8}{10}\selectfont
\begin{promptbox}
\begin{verbatim}
Conduct the triple-audit on this combined 2x2 grid image. Complete three tasks
one by one.

================================================================== TASK A -- T1
Cross-View Consistency (top-left vs top-right)
================================================================== Evaluate
whether T1 street-view (top-left quadrant) matches T1 satellite (top-right
quadrant). Analyze all four directions (North, East, South, West) relative to red
dot.

SCORING RUBRIC (0-10):
- 10: Perfect structural match. Differences purely
seasonal/color/lighting/quality.
- 7-9: Highly consistent. Very minor structural differences, macro environment
identical.
- 4-6: Noticeable structural changes in one or two directions (new/demolished
building, road altered).
- 1-3: Major mismatches across multiple directions.
- 0: [FATAL] Top-left street-view has blurred areas.

================================================================== TASK B -- T2
Cross-View Consistency (bottom-left vs bottom-right)
================================================================== Evaluate
whether T2 street-view (bottom-left) matches T2 satellite (bottom-right). SAME
criteria as Task A.

SCORING RUBRIC (0-10): same as Task A.
- 0: [FATAL] Bottom-left street-view has blurred areas.

================================================================== TASK C --
Satellite Temporal Consistency (top-right vs bottom-right)
================================================================== Compare T1
satellite (top-right) vs T2 satellite (bottom-right). Verify "NO-CHANGE" in
physical urban structures. Analyze all four directions relative to red dot.

SCORING RUBRIC (0-10):
- 10: Perfect structural match. Differences purely
seasonal/color/lighting/quality.
- 7-9: Highly consistent. Very minor structural differences (tiny shed, one tree
removed), macro environment identical.
- 4-6: Noticeable structural changes in one or two directions.
- 1-3: Major mismatches across multiple directions, large-scale changes.
- 0: Completely disjointed locations with massive contradictions.

================================================================== OUTPUT FORMAT
-- ALL THREE TASKS
================================================================== You MUST
output ONLY the three scores, nothing else, in this exact format:

[T1 Cross-View Score]: <Score 0-10>
[T2 Cross-View Score]: <Score 0-10>
[Satellite Consistency Score]: <Score 0-10>
\end{verbatim}
\end{promptbox}
\endgroup

\subparagraph{D3(b): Street-view temporal audit.}
\noindent\textbf{System prompt.}\par
\begingroup\fontsize{8}{10}\selectfont
\begin{promptbox}
\begin{verbatim}
You are an expert evaluator for a street-view dataset consistency audit.

Two 360-degree street-view panoramas of the SAME geographic location are provided
in a single composite image, one from an EARLIER period and one from a CURRENT
period:
- The TOP half is the HISTORICAL STREET VIEW (earlier period).
- The BOTTOM half is the CURRENT STREET VIEW (current period).

Your task is to verify whether this is a true "NO-CHANGE" location: the permanent
built environment should be structurally identical between the two captures.

Both panoramas are 360-degree images flattened onto a 2D canvas, pre-rotated so
the camera faces due North. Reference lines and labels are drawn on the images:
- CENTER (thick solid white line, yellow 'N') = North
- RIGHT-QUARTER (dashed gray line, 'E') = East
- LEFT-QUARTER (dashed gray line, 'W') = West
- BOTH side EDGES (dashed blue lines, 'S') = South

WHAT COUNTS AS A REAL CHANGE:
- BUILDINGS: any change to a building's static structure (built, demolished,
added/removed storeys, footprint change) OR its facade (cladding, paint
color/pattern, signage, windows, doors) counts as a REAL change.
- ROADS & PAVED SURFACES: changes to lane markings, crosswalks, road-surface
markings, curbs, or the paved ground surface count as a REAL change. IGNORE
surface differences caused purely by lighting, camera shooting style/exposure, or
water ponding on the surface.
- VEGETATION & TREES: changes to trees, hedges, bushes, or grass can be IGNORED —
vegetation is not a structural change.

ALWAYS IGNORE: lighting, weather, parked cars, pedestrians, shadows, seasonal
changes, and minor positional/scale/parallax shifts from different camera capture
points.
\end{verbatim}
\end{promptbox}
\endgroup

\noindent\textbf{User prompt.}\par
\begingroup\fontsize{8}{10}\selectfont
\begin{promptbox}
\begin{verbatim}
Evaluate whether the CURRENT street view (BOTTOM half) is genuinely UNCHANGED
compared to the HISTORICAL street view (TOP half), in terms of the permanent
built environment.

Consider all four directions (N/E/S/W) using the reference lines. Compare
buildings (structure and facade) and roads (surface and markings) between the two
halves. Ignore vegetation changes.

Score SV_T from 0 to 10:
- 10: Buildings and roads are structurally identical between the two street
views; no real change.
- 7-9: Essentially unchanged; only ignorable differences (lighting, weather,
seasonal, parked cars, vegetation).
- 4-6: Noticeable real changes in one or two directions (a building
structure/facade changed, or road markings/surface changed).
- 1-3: Major real changes across multiple directions.
- 0: Completely different scenes / massive redevelopment.

Output ONLY a single line — no reasoning, no explanation, no markdown. Just:
[SV_T]: <0-10>
\end{verbatim}
\end{promptbox}
\endgroup

\paragraph{D4: Satellite-grounded change descriptions (Gemini-3.1-Pro-Preview).}
\noindent\textbf{System prompt.}\par
\begingroup\fontsize{8}{10}\selectfont
\begin{promptbox}
\begin{verbatim}
You are an expert urban geographer and computer vision data annotator. Your task
is to align top-down satellite orthophotos with 360-degree street-level
panoramas. CRITICAL RULE: The street-view images have been pre-rotated so the
camera always faces due North. Vertical reference lines and text labels are
explicitly drawn on the street-view images to guide you: - Center = North (N):
Marked by a thick solid white line and yellow 'N' text. - Left-quarter = West
(W): Marked by a dashed gray line and 'W' text. - Right-quarter = East (E):
Marked by a dashed gray line and 'E' text. - Both left and right edges = South
(S): Marked by dashed blue lines and 'S' text. SPATIAL DIRECTION MAPPING: To
compare these two different perspectives, you must understand how their
directions align relative to the camera position (the RED DOT in the satellite
image): - North (N): The center of the street-view = looking Up (Top) from the
red dot. - East (E): The right-quarter of the street-view = looking Right from
the red dot.
- South (S): Both far-left and far-right edges wrap around = looking Down
(Bottom) from the red dot. - West (W): The left-quarter of the street-view =
looking Left from the red dot. Always describe scenes and map changes by strictly
following this directional mapping. On satellite images: N/S/W/E labels and the
north arrow are drawn in yellow with black outline. You must output exactly in
the requested bracketed formats.
\end{verbatim}
\end{promptbox}
\endgroup

\noindent\textbf{User prompt.}\par
\begingroup\fontsize{8}{10}\selectfont
\begin{promptbox}
\begin{verbatim}
You are given TWO annotated satellite images of the SAME geographic location:
Image 1 — HISTORICAL satellite image (older period) Image 2 — CURRENT satellite
image (recent period)

Each image has a RED DOT at the center (street-view capture point), a YELLOW
ARROW pointing due North (= up), and a small yellow text with black outline in
the top-left corner identifying it as "HISTORICAL satellite image" or "CURRENT
satellite image". Compare both images and identify permanent static structural
changes in the four cardinal directions relative to the red dot.

CRITICAL — GROUND-LEVEL PERSPECTIVE RULE: You must perform mental reconstruction
from the camera position (the RED DOT). Imagine you are standing at the red dot
looking outward in each direction at street level. Describe changes from this
ground-level viewpoint — what would actually be visible to a person standing
there. Do NOT report changes that are only visible from a bird's-eye view and
cannot be seen from the street (e.g., rooftop modifications, solar panel
additions).

WHAT QUALIFIES AS A STRUCTURAL CHANGE (MUST REPORT):
- Buildings: new construction, demolition, footprint expansion/shrinkage.
- Roads: significant changes to road surface material (e.g., dirt→paved) that are
clearly NOT caused by weather or lighting. Addition/removal/reconfiguration of
lane markings, signs, crosswalks, etc. Describe the pre-change structure in
sufficient detail (appearance, size, material, color).

WHAT TO IGNORE (MUST NOT REPORT):
- Seasonal and weather effects: wet vs. dry ground, leaf-on vs. leaf-off
vegetation, shadow differences.
- Vehicles and mobile objects: cars, trucks — regardless of how many or how
different they appear.
- Image quality differences: resolution, blur, color balance, brightness,
contrast between the two periods.

You MUST format your response EXACTLY with these four headers:

[North Direction (Up)]: State what existed historically and what exists now North
of the red dot(Describe the structure in sufficient detail (appearance, size,
material, color)). If unchanged, say "Unchanged: <description>".
[East Direction (Right)]: State what existed historically and what exists now
East of the red dot(Describe the structure in sufficient detail (appearance,
size, material, color)). If unchanged, say "Unchanged: <description>".
[South Direction (Bottom)]: State what existed historically and what exists now
South of the red dot(Describe the structure in sufficient detail (appearance,
size, material, color)). If unchanged, say "Unchanged: <description>".
[West Direction (Left)]: State what existed historically and what exists now West
of the red dot(Describe the structure in sufficient detail (appearance, size,
material, color)). If unchanged, say "Unchanged: <description>".

Keep each section factual, spatially precise, and under 100 words. Total response
under 400 words.
\end{verbatim}
\end{promptbox}
\endgroup

\paragraph{D5: Local editing instruction generation (Gemini-3.1-Pro-Preview).}
\noindent\textbf{System prompt.}\par
\begingroup\fontsize{8}{10}\selectfont
\begin{promptbox}
\begin{verbatim}
You are an expert urban geographer and computer vision data annotator. Your task
is to align top-down satellite orthophotos with 360-degree street-level
panoramas. CRITICAL RULE: The street-view images have been pre-rotated so the
camera always faces due North. Vertical reference lines and text labels are
explicitly drawn on the street-view images to guide you: - Center = North (N):
Marked by a thick solid white line and yellow 'N' text. - Left-quarter = West
(W): Marked by a dashed gray line and 'W' text. - Right-quarter = East (E):
Marked by a dashed gray line and 'E' text. - Both left and right edges = South
(S): Marked by dashed blue lines and 'S' text. SPATIAL DIRECTION MAPPING: To
compare these two different perspectives, you must understand how their
directions align relative to the camera position (the RED DOT in the satellite
image): - North (N): The center of the street-view = looking Up (Top) from the
red dot. - East (E): The right-quarter of the street-view = looking Right from
the red dot.
- South (S): Both far-left and far-right edges wrap around = looking Down
(Bottom) from the red dot. - West (W): The left-quarter of the street-view =
looking Left from the red dot. Always describe scenes and map changes by strictly
following this directional mapping. On satellite images: N/S/W/E labels and the
north arrow are drawn in yellow with black outline. You must output exactly in
the requested bracketed formats.
\end{verbatim}
\end{promptbox}
\endgroup

\noindent\textbf{User prompt.}\par
\begingroup\fontsize{8}{10}\selectfont
\begin{promptbox}
\begin{verbatim}
You are given:
1. A current T2 STREET-VIEW IMAGE (360° panorama, pre-rotated with camera facing
North, with N/S/W/E reference lines drawn).
2. A SATELLITE CHANGE ANALYSIS text describing structural differences between
historical and current satellite images.

SATELLITE CHANGE ANALYSIS: {layer1}

Your task: translate the satellite-confirmed changes into explicit 2D image
editing instructions. CRITICAL: The goal is to guide a Generative AI image
editing model to modify this exact 2D street-view image. You must perform mental
reconstruction — understanding spatial perspective, occlusion, and how each
viewing direction maps to a specific region on the 2D canvas.

STREET-VIEW DIRECTION → 2D CANVAS REGION MAPPING:
- North (N): CENTER of the image (thick solid white line + yellow 'N')
- East (E): RIGHT-QUARTER of the image (dashed gray line + 'E')
- South (S): BOTH the FAR-LEFT and FAR-RIGHT EDGES (dashed blue lines + 'S')
- West (W): LEFT-QUARTER of the image (dashed gray line + 'W')

You MUST output EXACTLY these four headers:

[Center (North View)]: <instruction>
[Right-quarter (East View)]: <instruction>
[Both side edges (South View)]: <instruction>
[Left-quarter (West View)]: <instruction>

Rules:
- Write "PRESERVE: keep all elements unchanged." if that direction has no
satellite-confirmed structural changes.
- Write "MODIFY: remove <current entity>, restore <historical entity>." if
changes occurred. Explicitly instruct how to handle occlusion, reveal
backgrounds, and adjust 3D perspective. Describe the structure in sufficient
detail (appearance, size, material, color).
- If a satellite change is occluded from ground view (e.g., rooftop modification
hidden by building facade), write PRESERVE.
- Be precise about spatial position (e.g., "on the left side of the center view",
"behind the foreground building on the right", "in the southeast direction", "on
the east side of the road").
- Keep each section under 150 words. Total response under 600 words.
\end{verbatim}
\end{promptbox}
\endgroup

\paragraph{D6: Instruction consistency validation (Gemini-3.1-Flash-Lite).}
Each criterion is evaluated in a separate call using the shared system message, its criterion-specific user prompt, and the editing instruction. All three scores must be at least 7; changed samples with all-`PRESERVE' instructions are removed.
The user message also appends \texttt{EDITING INSTRUCTION:} followed by the sample's instruction, then \texttt{COMPOSITE IMAGE (top half = CURRENT STREET VIEW, bottom half = HISTORICAL STREET VIEW):} and the composite image.
\noindent\textbf{Shared system prompt.}\par
\begingroup\fontsize{8}{10}\selectfont
\begin{promptbox}
\begin{verbatim}
You are an expert evaluator for a street-view image editing dataset.

A generative image-editing model is trained to edit the CURRENT STREET VIEW
panorama so that it becomes the HISTORICAL STREET VIEW panorama, guided by a
four-direction editing instruction written from satellite imagery.

- CURRENT STREET VIEW: a 360-degree panorama from the current period. This is the
model's INPUT; the model edits this image.
- HISTORICAL STREET VIEW: a 360-degree panorama from an earlier period. This is
the TARGET / ground-truth the model must reproduce.

The input image is a single COMPOSITE with the two panoramas stacked VERTICALLY,
each labeled at its top-left corner:
- The TOP half is the CURRENT STREET VIEW (model input).
- The BOTTOM half is the HISTORICAL STREET VIEW (target / ground-truth).

Both panoramas are 360-degree images flattened onto a 2D canvas, pre-rotated so
the camera faces due North. Reference lines and labels are drawn on the images:
- CENTER (thick solid white line, yellow 'N') = North
- RIGHT-QUARTER (dashed gray line, 'E') = East
- LEFT-QUARTER (dashed gray line, 'W') = West
- BOTH side EDGES (dashed blue lines, 'S') = South

The instruction has four blocks, one per direction, each marked MODIFY or
PRESERVE:
- "MODIFY: remove <current entity>, restore <historical entity>." — a structural
change happened in this direction.
- "PRESERVE: keep all elements unchanged." — no change happened in this
direction.

The two panoramas may be captured from slightly different camera positions; minor
positional/scale/parallax shifts between the two halves are EXPECTED and should
NOT be penalized. Ignore trivial differences such as lighting, weather, parked
cars, pedestrians, shadows, and seasonal vegetation. Only permanent structures
and ground elements matter.
\end{verbatim}
\end{promptbox}
\endgroup

\noindent\textbf{User prompt: Grounding.}\par
\begingroup\fontsize{8}{10}\selectfont
\begin{promptbox}
\begin{verbatim}
Evaluate ONE dimension: whether the editing instruction is EXECUTABLE on the
CURRENT street view (the TOP half of the composite).

Focus only on the MODIFY blocks. Evaluate EVERY MODIFY block independently. For
each MODIFY instruction, check whether the entity/region it tells the model to
remove is actually present and visible in the CURRENT street view at the stated
direction (use the N/E/S/W reference lines). If a MODIFY target appears in a
DIFFERENT direction than stated, treat it as NOT correctly grounded (major
location error).

The overall G is governed by the WORST-executable MODIFY block — a single
undetectable or misplaced target is a real defect even if other targets are
clear.

Score G from 0 to 10:
- 10: Every MODIFY target is clearly visible in the CURRENT street view at the
correct location; fully executable.
- 7-9: MODIFY targets are visible but with minor imprecision (slight location
drift or vague wording).
- 4-6: Some MODIFY targets are ambiguous, only partially visible, or in the wrong
direction.
- 1-3: Most MODIFY targets are not visible in the CURRENT street view (occluded,
absent, or hallucinated).
- 0: No MODIFY target can be found in the CURRENT street view at all.

Output ONLY a single line — no reasoning, no explanation, no markdown. Just:
[G]: <0-10>
\end{verbatim}
\end{promptbox}
\endgroup

\noindent\textbf{User prompt: Consistency.}\par
\begingroup\fontsize{8}{10}\selectfont
\begin{promptbox}
\begin{verbatim}
Evaluate ONE dimension: whether the HISTORICAL street view (the BOTTOM half of
the composite, the target / ground-truth) matches the MODIFY instructions.

For each MODIFY block, the instruction says "restore <historical entity>".
Evaluate EVERY MODIFY block independently. Check whether that historical entity
is actually present in the HISTORICAL street view at the stated direction, and
whether the HISTORICAL street view visibly differs from the CURRENT street view
in exactly the described way. If the historical entity appears in a DIFFERENT
direction than stated, treat it as a location mismatch.

The overall C is governed by the WORST-matching MODIFY block.

Score C from 0 to 10:
- 10: The HISTORICAL street view clearly shows every historical entity the
instruction says to restore; the instruction matches the ground truth.
- 7-9: It matches with minor discrepancies (extra or missing small objects).
- 4-6: It only partially matches; some described entities are absent, different,
or in the wrong direction in the HISTORICAL street view.
- 1-3: The HISTORICAL street view contradicts the instruction in most MODIFY
regions.
- 0: The HISTORICAL street view shows nothing resembling the historical state the
instruction describes.

Output ONLY a single line — no reasoning, no explanation, no markdown. Just:
[C]: <0-10>
\end{verbatim}
\end{promptbox}
\endgroup

\noindent\textbf{User prompt: Preservation.}\par
\begingroup\fontsize{8}{10}\selectfont
\begin{promptbox}
\begin{verbatim}
Evaluate ONE dimension: whether the regions marked PRESERVE in the instruction
are genuinely unchanged between the CURRENT and HISTORICAL street views.

For each PRESERVE block, compare the corresponding direction region between the
TOP half (CURRENT) and the BOTTOM half (HISTORICAL). Check whether permanent
structures in that region are structurally identical. Evaluate EVERY PRESERVE
block independently — any PRESERVE region with a real structural change drags the
overall P down.

Vegetation rule: a permanently removed or added tree/hedge counts as a structural
change; seasonal leaf/color changes do not.

Score P from 0 to 10:
- 10: All PRESERVE regions are structurally identical between the HISTORICAL and
CURRENT street views.
- 7-9: PRESERVE regions essentially unchanged; only ignorable differences
(seasonal, lighting, parked cars).
- 4-6: One or two PRESERVE regions contain noticeable structural changes the
instruction ignored.
- 1-3: Several PRESERVE regions contain clear structural changes the instruction
missed.
- 0: Major structural changes occur in PRESERVE regions; the instruction missed
them entirely.

If the instruction has NO PRESERVE blocks (all MODIFY), score P =
10.

Output ONLY a single line — no reasoning, no explanation, no markdown. Just:
[P]: <0-10>
\end{verbatim}
\end{promptbox}
\endgroup

\clearpage
